\documentclass[11pt]{article}
\usepackage[margin=0.88in]{geometry}
\usepackage{amsmath,amssymb,lmodern}
\usepackage[T1]{fontenc}
\usepackage{needspace}
\usepackage[colorlinks=true,linkcolor=blue,urlcolor=blue]{hyperref}
\setlength{\emergencystretch}{3em}
\setlength{\parskip}{0.35em}
\setcounter{secnumdepth}{0}
\providecommand{\tightlist}{\setlength{\itemsep}{0pt}\setlength{\parskip}{0pt}}
\title{Corner singularities obstruct eventual log-concavity of Robin ground states}
\author{Henry Zweiman}
\date{September 15, 2026}
\begin{document}
\maketitle
\section{Abstract}\label{abstract}

Let \(P\) be a bounded convex polygon that is neither tangential nor a
rectangle. We prove that the positive Robin parameters for which the
first eigenfunction is log-concave form a locally finite set. In fact,
at one fixed obtuse vertex, the eigenfunction and its logarithm fail to
be semiconcave on every corner neighborhood for every parameter outside
the zero set of a nontrivial real-analytic function. The proof continues
a corner singularity of the first variation at the Neumann limit. An
explicit angular projection gives an analytic singular coefficient
without requiring analytic dependence of a full corner expansion. This
disproves the eventual log-concavity conjecture of Andrews, Clutterbuck
and Hauer for general bounded convex domains. A single nonrhombic,
nonrectangular parallelogram gives a counterexample, and products give
examples in every dimension at least two. An interval-lifting argument
further gives fixed convex prisms in every dimension at least three
whose ground states have nonconvex superlevel sets outside the same
discrete parameter set. A weighted Neumann argument also proves
transverse nonconcavity of homogeneous modes on a specified class of
higher-dimensional cones, supplying a conditional small-parameter Robin
consequence. The result is compatible with eventual concavity on smooth
uniformly convex domains.

\textbf{Keywords:} Robin eigenfunction; log-concavity; polygonal domain;
corner singularity; analytic perturbation; quasiconcavity.

\textbf{2020 Mathematics Subject Classification:} 35J25, 35B65, 35P05,
47A55.

\section{1. Introduction}\label{introduction}

For a bounded convex polygon \(P\subset\mathbb R^2\) with nonempty
interior, let \(u_\alpha\) be the positive first Robin eigenfunction:

\[
-\Delta u_\alpha=\lambda_\alpha u_\alpha\quad\text{in }P,
\qquad \partial_\nu u_\alpha+\alpha u_\alpha=0
\quad\text{on }\partial P.
\tag{1.1}
\]

The boundary condition is understood weakly; on the interior of an edge
it is classical. We fix the normalization

\[
\int_P u_\alpha^2\,dx=|P|.
\tag{1.2}
\]

In particular, \(u_0=1\) and \(\lambda_0=0\). Log-concavity means
concavity of \(\log u_\alpha\) in the open convex set \(P\).

Andrews, Clutterbuck and Hauer {[}ACH{]} proved failure of log-concavity
for sufficiently small positive Robin parameters on convex polyhedra
outside a geometrically distinguished class. Their final Section 10,
Conjecture 1, asks whether every fixed bounded convex domain has a
threshold above which its Robin ground state is log-concave. Crasta and
Fragalà {[}CF{]} established eventual log-concavity on sufficiently
smooth uniformly convex domains. Ye and Zhang {[}YZ{]} recently reduced
the boundary assumption to \(C^{3,1}\) in every dimension. These
theorems leave open the conjecture on convex domains with corners.

A polygon is \emph{tangential} if it contains a disk tangent to every
edge. This is the planar instance of a circumsolid in {[}ACH{]}. In
dimension two, products of circumsolids are precisely tangential
polygons and rectangles. We take polygonal descriptions with no
redundant collinear vertices, so every interior angle belongs to
\((0,\pi)\).

\Needspace{6\baselineskip}\textbf{Theorem 1.1.} Let \(P\) be a bounded convex polygon that is
neither tangential nor a rectangle. There exist an obtuse vertex \(v_*\)
and a real-analytic function \(c:\mathbb R\to\mathbb R\) with

\[
c(0)=0,\qquad c'(0)\ne0,
\tag{1.3}
\]

such that, whenever \(\alpha>0\) and \(c(\alpha)\ne0\), neither
\(u_\alpha\) nor \(\log u_\alpha\) is semiconcave on any sufficiently
small neighborhood \(P\cap B_r(v_*)\). Consequently

\[
\{\alpha>0:u_\alpha\text{ is log-concave in }P\}
\subseteq\{\alpha>0:c(\alpha)=0\}.
\tag{1.4}
\]

The set on the left is locally finite, and it is empty for all
sufficiently small positive parameters. In particular, for every \(A>0\)
there is an \(\alpha>A\) for which \(u_\alpha\) is not log-concave.

Here semiconcavity on a convex set means that \(f(x)-M|x|^2/2\) is
concave for some finite constant \(M\). The assertion near \(v_*\)
refers to each sufficiently small radius, with no uniformity in
\(\alpha\) required. In fact, the exceptional zero set is finite on
every bounded parameter interval, including intervals meeting zero.

\Needspace{6\baselineskip}\textbf{Corollary 1.2.} The polygon with consecutive vertices

\[
(0,0),\quad(2,0),\quad(3,1),\quad(1,1)
\tag{1.5}
\]

does not have an eventual log-concavity threshold. There are bounded
convex polyhedral counterexamples in every dimension \(d\ge2\).

The new point is the continuation of a regularity obstruction through
the entire finite Robin parameter axis. The first-variation
classification in {[}ACH{]} supplies the initial nonzero corner mode. We
extract its continuation by a bounded functional on a fixed Sobolev
space, then use the identity theorem for real-analytic functions. Thus
the argument does not extrapolate a small-parameter asymptotic estimate
to large parameters.

We do not determine whether log-concavity occurs at any of the
exceptional parameters. Nor do we prove a classification for tangential
polygons. Section 6 gives nonconvex superlevel sets on fixed prisms in
every dimension at least three for all parameters outside the same
discrete set. It does not resolve {[}ACH, Section 10, Conjecture 2{]},
which concerns every convex polyhedron that is not a product of
circumsolids for sufficiently small parameters, or the sharp Robin
fundamental-gap problem.

\section{2. Analytic dependence and local
regularity}\label{analytic-dependence-and-local-regularity}

We first record precisely the regularity needed for the coefficient
construction. Write

\[
\Gamma_\omega=\{(r\cos\theta,r\sin\theta):r>0,\ 0<\theta<\omega\},
\qquad S_R=\Gamma_\omega\cap B_R(0),
\qquad \beta=\frac\pi\omega>1.
\tag{2.1}
\]

\Needspace{6\baselineskip}\textbf{Lemma 2.1 (sector regularity).} Suppose
\(z\in H^1(\Gamma_\omega)\) is supported in a fixed closed ball, has
homogeneous weak Neumann data on the rays, and
\(\Delta z\in L^p(\Gamma_\omega)\), where \(2\le p<\infty\). If

\[
2-\frac2p<\beta,
\tag{2.2}
\]

then \(z\in W^{2,p}(\Gamma_\omega)\). For a fixed support radius the
inclusion is continuous for the graph norm
\(\|z\|_{H^1}+\|\Delta z\|_{L^p}\).

\emph{Proof.} This is a direct planar consequence of
Dauge\textquotesingle s wedge theorem {[}D, Theorem 8.1{]}, which
applies to compactly supported weak Neumann solutions on
\(\mathbb R\times\Gamma_\omega\) when \(0<2-2/p<\pi/\omega\). To check
the reduction, choose a nonzero \(\eta\in C_c^\infty(\mathbb R)\) and
set \(Z(t,x)=\eta(t)z(x)\). Then \(Z\in H^1\) has compact support and

\[
\Delta_{t,x}Z=\eta\Delta_xz+\eta''z\in L^p.
\tag{2.3}
\]

Here \(z\in L^p\) follows from the planar \(H^1\) embedding on a bounded
sector containing its support. The boundary data of \(Z\) are
homogeneous Neumann. The cited theorem gives \(Z\in W^{2,p}\), hence
\(z\in W^{2,p}\), for example by integration against a smooth function
\(\xi(t)\) with \(\int\xi\eta=1\).

For continuity, consider the pairs \((z,f)\in H^1\times L^p\) with the
fixed support restriction and

\[
\int_{\Gamma_\omega}\nabla z\cdot\nabla\varphi
=-\int_{\Gamma_\omega}f\varphi
\quad\text{for all }\varphi\in C_c^\infty(\overline{\Gamma_\omega}).
\tag{2.4}
\]

They form a closed subspace of \(H^1\times L^p\). The map
\((z,f)\mapsto z\) into \(W^{2,p}\) has closed graph, since both limits
agree as distributions. The closed graph theorem proves the estimate. \(\square\)

\Needspace{6\baselineskip}\textbf{Lemma 2.2 (analytic eigenpair).} With normalization (1.2), the
maps \(\alpha\mapsto\lambda_\alpha\) and \(\alpha\mapsto u_\alpha\) are
real analytic on \(\mathbb R\), the latter with values in \(H^1(P)\).

\emph{Proof.} The Robin form

\[
a_\alpha(f,g)=\int_P\nabla f\cdot\nabla g
+\alpha\int_{\partial P}fg
\tag{2.5}
\]

has fixed domain \(H^1(P)\). The trace map is bounded, and a
sufficiently large scalar shift makes the form coercive locally
uniformly in \(\alpha\). The first eigenvalue is simple and its
eigenfunction can be chosen positive; see {[}ACH, Proposition 3.1{]}.

For completeness, analytic dependence follows from the analytic implicit
function theorem applied to

\[
(\alpha,u,\lambda)\longmapsto
\left(a_\alpha(u,\cdot)-\lambda(u,\cdot)_{L^2(P)},
\ (u,u)_{L^2(P)}-|P|\right)
\tag{2.6}
\]

with values in \((H^1(P))^*\times\mathbb R\). Its derivative in
\((u,\lambda)\) sends \((h,s)\) to

\[
\left(B_\alpha h-s(u_\alpha,\cdot)_{L^2(P)},
\ 2(u_\alpha,h)_{L^2(P)}\right),
\quad B_\alpha=a_\alpha-\lambda_\alpha(\cdot,\cdot)_{L^2(P)}.
\tag{2.7}
\]

The operator \(B_\alpha:H^1\to(H^1)^*\) is Fredholm of index zero, by a
coercive shift and compactness of the \(L^2\) embedding. Its kernel is
the span of \(u_\alpha\), and its range is the annihilator of
\(u_\alpha\). Pairing the first component with \(u_\alpha\) determines
\(s\), the second component determines the component of \(h\) along
\(u_\alpha\), and \(B_\alpha\) is invertible on the complementary
subspace. Thus (2.7) is an isomorphism. The local analytic branches
remain the first eigenpair by the spectral gap. Positivity and (1.2)
make them agree on overlaps, giving the asserted real-analytic maps. \(\square\)

At a vertex translated to the origin, let \(\nu_1,\nu_2\) be the two
outward unit normals. They are linearly independent. Let \(\gamma\) be
the unique vector satisfying

\[
\nu_1\cdot\gamma=\nu_2\cdot\gamma=-1.
\tag{2.8}
\]

Define the local gauge

\[
w_\alpha(x)=e^{-\alpha\gamma\cdot x}u_\alpha(x).
\tag{2.9}
\]

It satisfies homogeneous Neumann conditions on both rays and

\[
\Delta w_\alpha=F_\alpha,
\qquad
F_\alpha=-2\alpha\gamma\cdot\nabla w_\alpha
-(\lambda_\alpha+\alpha^2|\gamma|^2)w_\alpha.
\tag{2.10}
\]

\Needspace{6\baselineskip}\textbf{Lemma 2.3 (regularity of the gauge).} On a sufficiently small
fixed sector, \(w_\alpha\) depends real analytically on \(\alpha\) in
\(H^2\). Consequently \(F_\alpha\) is real analytic in \(L^p\) for every
finite \(p\). For each fixed \(\alpha\), \(w_\alpha\) is
\(C^{1,\delta}\) up to the vertex for some \(\delta>0\), and

\[
\nabla w_\alpha(0)=0.
\tag{2.11}
\]

\emph{Proof.} Choose a smooth radial cutoff \(\chi\) supported strictly
inside the radius on which \(P\) agrees with the sector, and equal to
one on a smaller sector. Extend \(z_\alpha=\chi w_\alpha\) by zero to
\(\Gamma_\omega\). Its ray normal derivative is zero. Lemma 2.2 and
multiplication by the exponential show that \(z_\alpha\) is analytic in
\(H^1\). Moreover,

\[
\Delta z_\alpha
=\chi F_\alpha+2\nabla\chi\cdot\nabla w_\alpha
+(\Delta\chi)w_\alpha
\tag{2.12}
\]

is analytic in \(L^2\). The pairs \((z_\alpha,\Delta z_\alpha)\) are
therefore analytic in the closed graph space in Lemma 2.1. Since
\(1<\beta\), that lemma with \(p=2\) makes \(z_\alpha\) analytic in
\(H^2\). Planar Sobolev embedding gives analytic dependence of
\(w_\alpha\) and \(\nabla w_\alpha\) in every finite \(L^p\) on a
smaller sector, and hence the claim about \(F_\alpha\).

Choose \(p>2\) sufficiently close to two that (2.2) holds. A second,
nested radial cutoff and (2.12), now with \(L^p\) data, give
\(w_\alpha\in W^{2,p}\) locally. Morrey embedding gives \(C^{1,\delta}\)
regularity, with \(0<\delta<1-2/p\). The two Neumann conditions then
hold at the vertex by continuity. The independent normals force (2.11).
\(\square\)

\Needspace{6\baselineskip}\textbf{Lemma 2.4 (positivity at the vertex).} If \(\alpha>0\), then
\(w_\alpha(0)>0\).

\emph{Proof.} On a sufficiently small sector put
\(\rho(x)=e^{2\alpha\gamma\cdot x}\) and
\(k=\lambda_\alpha+\alpha^2|\gamma|^2>0\). Equation (2.10) becomes

\[
-\operatorname{div}(\rho\nabla w_\alpha)=k\rho w_\alpha.
\tag{2.13}
\]

The ground state is positive in the interior. It is also positive at
every point in the relative interior of either edge: otherwise the Hopf
boundary lemma at that smooth boundary point would contradict the Robin
condition. Thus \(w_\alpha\) has a positive minimum \(m\) on the closed
outer circular arc of a small sector. Regularity from Lemma 2.3 gives
continuity on its closure.

Use \(\varphi=(m-w_\alpha)_+\) as a test function in (2.13), with zero
trace on the outer arc and no restriction on the rays. The ray fluxes
vanish, and \(w_\alpha\ge0\). Hence

\[
-\int_{S_R}\rho|\nabla\varphi|^2
=k\int_{S_R}\rho w_\alpha\varphi\ge0.
\tag{2.14}
\]

It follows that \(\varphi=0\), since its outer-arc trace is zero.
Therefore \(w_\alpha\ge m\) on the whole sector, including its vertex. \(\square\)

\section{3. A bounded functional for the corner
coefficient}\label{a-bounded-functional-for-the-corner-coefficient}

Fix an obtuse sector, so \(1<\beta<2\), and a radius \(R\) lying inside
the region of Lemma 2.3. Set

\[
\begin{aligned}
a_\alpha(r)&=\frac2\omega\int_0^\omega
w_\alpha(r,\theta)\cos(\beta\theta)\,d\theta,\\
b_\alpha(r)&=\frac2\omega\int_0^\omega
F_\alpha(r,\theta)\cos(\beta\theta)\,d\theta.
\end{aligned}
\tag{3.1}
\]

\Needspace{6\baselineskip}\textbf{Lemma 3.1 (analytic coefficient).} The limit

\[
c(\alpha)=\lim_{r\downarrow0}r^{-\beta}a_\alpha(r)
\tag{3.2}
\]

exists for every real \(\alpha\) and is a real-analytic function of
\(\alpha\). More precisely,

\[
\begin{aligned}
c(\alpha)={}&R^{-\beta}a_\alpha(R)\\
&-\frac1{2\beta}\int_0^R
\big(s^{1-\beta}-R^{-2\beta}s^{1+\beta}\big)b_\alpha(s)\,ds.
\end{aligned}
\tag{3.3}
\]

\emph{Proof.} Angular projection of (2.10), using the homogeneous
Neumann conditions, gives

\[
a_\alpha''+\frac1r a_\alpha'-\frac{\beta^2}{r^2}a_\alpha=b_\alpha.
\tag{3.4}
\]

This equation first holds distributionally away from zero; the functions
are smooth there. Choose a finite exponent \(p\) so large that

\[
p>\frac2{2-\beta},\qquad \kappa:=2-\frac2p>\beta.
\tag{3.5}
\]

Lemma 2.3 permits this choice for \(F_\alpha\). Variation of constants
in (3.4) gives the regular particular solution

\[
A_\alpha(r)=\frac1{2\beta}\left[
r^\beta\int_0^r s^{1-\beta}b_\alpha(s)\,ds
-r^{-\beta}\int_0^r s^{1+\beta}b_\alpha(s)\,ds\right].
\tag{3.6}
\]

Hölder\textquotesingle s inequality, with planar area measure
\(s\,ds\,d\theta\), gives

\[
|A_\alpha(r)|+r|A_\alpha'(r)|
\le C r^\kappa\|F_\alpha\|_{L^p(S_r)}.
\tag{3.7}
\]

For the first integral the relevant weight is \(s^{-\beta}\) relative to
area measure. Its \(L^{p'}\) norm is finite exactly when \(\beta p'<2\),
which is equivalent to (3.5). The second integral is less singular. The
homogeneous solutions of (3.4) are \(r^\beta\) and \(r^{-\beta}\). The
latter is excluded by \(w_\alpha\in H^1(S_R)\), since its angular
projection would have infinite \(H^1\) energy. Thus

\[
a_\alpha(r)=c(\alpha)r^\beta+A_\alpha(r).
\tag{3.8}
\]

Equations (3.7) and (3.5) prove (3.2). Evaluating (3.8) at \(R\) proves
(3.3).

The trace functional \(w\mapsto a(R)\) is bounded on \(H^2\) on a sector
extending past \(R\). The integral in (3.3) is a bounded linear
functional of \(F\in L^p(S_R)\) by the same weighted Hölder estimate.
Lemma 2.3 therefore proves analyticity of (3.3). The value is
independent of the auxiliary radius because it equals the limit (3.2). \(\square\)

Two different exponents have been used for different purposes. An
exponent just above two yields \(C^1\) regularity in Lemma 2.3. A much
larger exponent is allowed for the forcing \(F_\alpha\) in Lemma 3.1. No
\(W^{2,p}\) regularity of \(w_\alpha\) is asserted for that larger
exponent; generally it is false when \(c(\alpha)\ne0\).

\Needspace{6\baselineskip}\textbf{Lemma 3.2 (semiconcavity annihilates the coefficient).} For
\(\alpha>0\), if either \(u_\alpha\) or \(\log u_\alpha\) is semiconcave
on a neighborhood of the vertex, then \(c(\alpha)=0\).

\emph{Proof.} By Lemmas 2.3 and 2.4, \(u_\alpha\), \(w_\alpha\) and
their logarithms are \(C^1\) on a smaller closed sector, and the
positive functions are bounded away from zero there. If \(u_\alpha\) is
semiconcave, its interior Hessian has an upper bound, while

\[
\Delta u_\alpha=-\lambda_\alpha u_\alpha
\tag{3.9}
\]

is bounded. In two dimensions a symmetric matrix with bounded trace and
eigenvalues bounded above has eigenvalues bounded below as well.
Therefore \(D^2u_\alpha\) is bounded in the interior of that sector.
Integration along segments in the convex sector gives
\(u_\alpha\in C^{1,1}\) up to the vertex, and multiplication by the
smooth gauge gives \(w_\alpha\in C^{1,1}\).

If \(\log u_\alpha\) is semiconcave, put
\(q=\log w_\alpha=\log u_\alpha-\alpha\gamma\cdot x\). Subtracting a
linear function preserves semiconcavity. Equation (2.10) gives

\[
\Delta q=-|\nabla q|^2-2\alpha\gamma\cdot\nabla q
-(\lambda_\alpha+\alpha^2|\gamma|^2).
\tag{3.10}
\]

Its right hand side is bounded. The same Hessian argument gives
\(q\in C^{1,1}\) and hence \(w_\alpha\in C^{1,1}\).

In either case, (2.11) yields

\[
w_\alpha(x)=w_\alpha(0)+O(|x|^2).
\tag{3.11}
\]

The angular integral of \(\cos(\beta\theta)\) is zero, so (3.1) gives
\(a_\alpha(r)=O(r^2)\). Since \(\beta<2\), (3.2) forces \(c(\alpha)=0\).
The same argument works on every smaller corner neighborhood. \(\square\)

\section{4. The initial nonzero mode}\label{the-initial-nonzero-mode}

Let

\[
v=\left.\frac{d u_\alpha}{d\alpha}\right|_{\alpha=0},
\qquad \mu=\frac{|\partial P|}{|P|}.
\tag{4.1}
\]

Differentiation of the weak eigenvalue equation and normalization gives

\[
\Delta v=-\mu\quad\text{in }P,
\qquad \partial_\nu v=-1\quad\text{on }\partial P,
\qquad \int_Pv=0.
\tag{4.2}
\]

We use the following established rigidity result, with regularity on the
closure understood: by {[}ACH, Corollary 8.3{]}, a solution of (4.2) is
\(C^2(\overline P)\) if and only if \(P\) is a product of circumsolids.
Thus it is not \(C^2(\overline P)\) under the hypotheses of Theorem 1.1.

\Needspace{6\baselineskip}\textbf{Lemma 4.1 (nonzero first variation).} If \(P\) is neither
tangential nor a rectangle, there is an obtuse vertex for which the
coefficient in Lemma 3.1 satisfies \(c(0)=0\) and \(c'(0)\ne0\).

\emph{Proof.} At any vertex, translate the vertex to zero and choose
\(\gamma\) as in (2.8). The function

\[
h(x)=v(x)+\frac\mu4|x|^2-\gamma\cdot x
\tag{4.3}
\]

is harmonic with homogeneous Neumann data on the two rays. Indeed, the
normal derivative of \(|x|^2\) vanishes on either ray. Separation of
variables gives

\[
h(r,\theta)=d_0+\sum_{j\ge1}d_j r^{j\pi/\omega}
\cos\left(\frac{j\pi\theta}{\omega}\right).
\tag{4.4}
\]

This expansion also follows from {[}ACH, Proposition 5.5{]}. It can be
obtained directly by expanding the trace on an inner arc in the Neumann
cosine basis. The \(H^1\) condition excludes negative powers and the
logarithmic zero mode. On every strictly smaller sector the series and
its derivatives converge away from zero. Moreover, the part with
exponents strictly greater than two has second derivatives tending to
zero at the vertex: square summability of the trace coefficients and the
geometric factor \((r/R)^{j\pi/\omega}\) justify termwise bounds on
smaller radii.

If \(\omega<\pi/2\), every nonconstant exponent exceeds two. If
\(\omega=\pi/2\), the first term is a harmonic quadratic polynomial, and
all subsequent exponents exceed two. If \(\pi/2<\omega<\pi\), only the
first exponent \(\beta=\pi/\omega\) lies below two, and the next is
\(2\beta>2\). Therefore failure of \(C^2\) regularity at a vertex is
possible only at an obtuse vertex with \(d_1\ne0\).

The solution \(v\) is smooth in the interior and up to the relative
interiors of edges, by reflection after subtracting a local particular
solution with the given normal derivative. If every vertex had a \(C^2\)
expansion, these local statements would give \(v\in C^2(\overline P)\),
contrary to the rigidity result above. Fix an obtuse vertex with
\(d_1\ne0\).

At this vertex \(w_0=1\), so \(c(0)=0\). Differentiating (2.9) and
(2.10) at zero gives

\[
\dot w_0=v-\gamma\cdot x,
\qquad \dot F_0=-\mu.
\tag{4.5}
\]

The constant forcing has zero first cosine projection. The radial
quadratic term in (4.3) also has zero first cosine projection. It
follows from (3.3), which may be differentiated by Lemma 3.1, and (4.4)
that

\[
c'(0)=d_1\ne0.
\tag{4.6}
\]

This is the claimed vertex and coefficient. \(\square\)

\emph{Proof of Theorem 1.1.} Choose the vertex from Lemma 4.1 and its
analytic coefficient from Lemma 3.1. Since \(c'(0)\ne0\), this
real-analytic function is not identically zero. Its zero set is locally
finite on \(\mathbb R\), by the identity theorem, and zero is an
isolated simple zero. Lemma 3.2 shows that semiconcavity of either
\(u_\alpha\) or its logarithm on any sufficiently small corner
neighborhood forces \(c(\alpha)=0\). This proves the local assertion and
(1.4). A locally finite subset cannot contain any interval \((A,A+1)\),
so there are arbitrarily large parameters outside the exceptional set. \(\square\)

\section{5. Explicit domains and the Dirichlet
limit}\label{explicit-domains-and-the-dirichlet-limit}

\emph{Proof of Corollary 1.2.} The polygon (1.5) is a parallelogram with
adjacent side lengths \(2\) and \(\sqrt2\), and angles \(\pi/4\) and
\(3\pi/4\). It is not a rectangle. A tangential quadrilateral has equal
sums of opposite side lengths: the two tangent segments from each vertex
have equal lengths, and summing them gives the identity. A tangential
parallelogram must therefore have equal adjacent side lengths. Our
polygon is not tangential, so Theorem 1.1 applies. At its obtuse
vertices the relevant exponent is \(\beta=4/3\).

For \(d>2\), use

\[
\Omega_d=P\times(0,1)^{d-2}.
\tag{5.1}
\]

The positive product of the first Robin eigenfunctions on the factors
satisfies the Robin condition with the same parameter on every face. Its
eigenvalue is the sum of the first eigenvalues of the factors, and
positivity identifies it with the ground state. Restriction to a slice
with fixed interior coordinates in the interval factors is a positive
constant multiple of \(u_\alpha\) on \(P\). Log-concavity of the product
would imply log-concavity on this slice. Thus every parameter excluded
for \(P\) is excluded for \(\Omega_d\). \(\square\)

The contrast with the Dirichlet limit concerns regularity at the
corners. A nonzero coefficient in (3.2) is compatible with convergence
in weaker norms and with arbitrarily small amplitude as \(\alpha\)
increases. For each fixed finite parameter with \(c(\alpha)\ne0\),
however, the mode \(r^\beta\) with \(\beta<2\) prevents \(C^{1,1}\)
regularity of the gauge at the vertex. Semiconcavity would force that
regularity through the bounded-trace argument of Lemma 3.2. No uniform
\(C^2\) convergence near the polygon\textquotesingle s corners is
available from a Dirichlet limit alone.

The results of {[}CF{]} and {[}YZ{]} assume smooth uniformly convex
domains. Our polygonal examples satisfy neither boundary hypothesis.
Approximating a polygon by smooth domains for a fixed parameter does not
produce a single smooth domain with the discrete-exception property
proved here, since the approximation can depend on the parameter. The
remaining smooth-boundary questions must be treated separately.

The method suggests two further problems. One is to determine the
exceptional parameters, or to show their absence for particular
nonsymmetric polygons. Another is to understand eventual log-concavity
within the tangential class, where the first-variation singular
coefficient vanishes and the present initial-mode argument gives no
obstruction. These questions are distinct from the conjecture for all
convex domains, which Corollary 1.2 answers negatively.

\section{6. Nonconvex superlevel sets on fixed
prisms}\label{nonconvex-superlevel-sets-on-fixed-prisms}

A positive function \(U\) on a convex domain is \emph{quasiconcave} if
every superlevel set \(\{U>t\}\) is convex. For a smooth positive
function, a strictly positive Hessian in a direction tangent to a level
set of \(\log U\) rules out quasiconcavity. We give the elementary
argument in the following lemma.

\Needspace{6\baselineskip}\textbf{Lemma 6.1 (interval lifting).} Let \(D\subset\mathbb R^m\) be an
open convex set, let \(f\in C^2(D)\) have bounded gradient, and suppose
\(f\) is not semiconcave on \(D\). Let \(I\) be an open interval and
\(h\in C^2(I)\) satisfy \(h'(t_0)\ne0\) at some \(t_0\in I\). Then

\[
U(x,t)=\exp(f(x)+h(t))
\tag{6.1}
\]

has a nonconvex superlevel set in \(D\times I\).

\emph{Proof.} If \(D^2 f(x)[v,v]\le M\) for every \(x\in D\) and every
unit vector \(v\), integration along segments shows that \(f-M|x|^2/2\)
is concave. Therefore, failure of semiconcavity gives points
\(x_j\in D\) and unit vectors \(v_j\) such that

\[
D^2f(x_j)[v_j,v_j]\longrightarrow+\infty.
\tag{6.2}
\]

Put \(F=\log U\) and

\[
s_j=-\frac{\nabla f(x_j)\cdot v_j}{h'(t_0)},
\qquad V_j=(v_j,s_j).
\tag{6.3}
\]

The numbers \(s_j\) are bounded. At \(z_j=(x_j,t_0)\) we have

\[
\nabla F(z_j)\cdot V_j=0,\qquad
D^2F(z_j)[V_j,V_j]
=D^2f(x_j)[v_j,v_j]+h''(t_0)s_j^2\longrightarrow+\infty.
\tag{6.4}
\]

Fix \(j\) for which the last expression is positive. Since \(z_j\) is
interior, Taylor\textquotesingle s theorem gives, for sufficiently small
\(\varepsilon>0\),

\[
F(z_j\pm\varepsilon V_j)>F(z_j).
\tag{6.5}
\]

Choose \(b\) strictly between \(F(z_j)\) and the smaller endpoint value.
Both endpoints belong to \(\{U>e^b\}\), while their midpoint does not.
This proves the assertion. \(\square\)

\Needspace{6\baselineskip}\textbf{Theorem 6.2 (persistent failure of quasiconcavity on prisms).}
Let \(P\) satisfy the hypotheses of Theorem 1.1, and let \(c\) be its
analytic corner coefficient. For every integer \(d\ge3\), the first
Robin eigenfunction on the fixed convex polyhedron

\[
\Omega_d=P\times(0,1)^{d-2}
\tag{6.6}
\]

has a nonconvex superlevel set whenever \(\alpha>0\) and
\(c(\alpha)\ne0\). Consequently, the parameters at which this ground
state is quasiconcave form a locally finite set, and there is no
eventual quasiconcavity threshold for \(\Omega_d\).

\emph{Proof.} Fix such an \(\alpha\). On a sufficiently small convex
corner neighborhood \(D=P\cap B_r(v_*)\), the function
\(f=\log u_\alpha\) is smooth in the interior and has bounded gradient,
by Lemmas 2.3 and 2.4. It is not semiconcave there, by Theorem 1.1.

The positive first Robin eigenfunction on \((0,1)\) can be written

\[
g_\alpha(t)=\cos(k(t-1/2)),\qquad
k\tan(k/2)=\alpha,\quad 0<k<\pi.
\tag{6.7}
\]

The equation for \(k\) has a unique solution because its left side is
strictly increasing from zero to infinity. Direct differentiation
verifies the Robin condition at both endpoints; positivity identifies
the eigenfunction as the ground state. For \(h=\log g_\alpha\) and
\(t_0=1/4\),

\[
h'(t_0)=k\tan(k/4)>0,
\qquad h''(t_0)=-k^2\sec^2(k/4).
\tag{6.8}
\]

Lemma 6.1 shows that \(u_\alpha(x)g_\alpha(t)\) has a nonconvex
superlevel set already in \(D\times(0,1)\). This product is the ground
state on \(P\times(0,1)\), by the product argument in Section 5. For
\(d>3\), fix the remaining interval coordinates at interior points. The
resulting slice is a positive constant multiple of the same product, so
the same two endpoints and midpoint witness nonconvexity. The assertions
about parameters follow from the locally finite zero set of \(c\). \(\square\)

In particular, the explicit parallelogram (1.5) gives a fixed prism in
each dimension \(d\ge3\) with nonconvex ground-state superlevel sets for
arbitrarily large Robin parameters. The small-parameter conclusion for
these product domains follows already from the planar result in
{[}ACH{]} and restriction to a slice. The additional conclusion here is
persistence outside a discrete set over the entire positive parameter
axis. No novelty is asserted for the elementary level-set criterion or
separation of variables.

This extension answers the eventual-quasiconcavity question for these
fixed prisms negatively. It leaves {[}ACH, Section 10, Conjecture 2{]}
unresolved on general polyhedra. The proof uses both bounded gradient
and unbounded positive Hessian near a polygonal corner; failure of
log-concavity alone would not imply the interval-lifting conclusion. No
single smooth uniformly convex domain with arbitrarily large bad
parameters is constructed here.

\section{7. A weighted obstruction on higher-dimensional
cones}\label{a-weighted-obstruction-on-higher-dimensional-cones}

The unresolved local condition in {[}ACH, Lemma 9.5 and Remark 9.6{]}
concerns the restriction of a homogeneous harmonic Neumann mode to a
transverse affine slice. The following result verifies that condition
for a geometric class of cones. It imposes no symmetry on the mode.

\Needspace{6\baselineskip}\textbf{Theorem 7.1 (transverse nonconcavity).} Let
\(\Gamma\subset\mathbb R^d\), \(d\ge3\), be an open full-dimensional
convex polyhedral cone with outward unit face normals \(\nu_i\). Suppose
a unit vector \(e\) and a number \(\eta>0\) satisfy

\[
e\cdot x\ge\eta|x|\quad(x\in\overline\Gamma),
\qquad 0<k_i:=-\nu_i\cdot e<\frac1{\sqrt2}
\quad\text{for every face}.
\tag{7.1}
\]

Let \(\psi\) be a nonzero homogeneous weak Neumann harmonic function of
degree \(\beta\in(1,2)\), with \(\psi\in H^1(\Gamma\cap B_R)\) for every
finite \(R\). Then \(\psi\) is not concave on the affine slice

\[
K=\{x\in\Gamma:e\cdot x=1\}.
\tag{7.2}
\]

\emph{Proof.} We first justify the derivatives used below. A smooth
radial cutoff \(\chi\) preserves the zero Neumann condition on the cone
faces. On a sufficiently large convex truncation \(\Gamma\cap B_R\), the
function \(z=\chi\psi\) vanishes near the outer sphere and has

\[
\Delta z=2\nabla\chi\cdot\nabla\psi+(\Delta\chi)\psi\in L^2.
\tag{7.3}
\]

The standard homogeneous Neumann regularity theorem on bounded convex
domains gives \(z\in H^2\). The theorem and its convex approximation
argument are recalled in {[}T, Section 4, following Theorem 4.2{]}. Thus
\(\psi\) is locally \(H^2\) up to the cone boundary, including its
vertex. In particular, \(h=\partial_e\psi\) is locally \(H^1\) and
harmonic in the interior. Even reflection gives smoothness up to every
relative face interior.

Suppose that the restriction to \(K\) is concave. The first condition in
(7.1) puts every nonzero cone point at positive height, so homogeneity
gives concavity on every parallel slice. Put \(H=D^2\psi\). Its
restriction to \(e^\perp\) is negative semidefinite, and harmonicity
gives

\[
q:=H[e,e]=-\operatorname{tr}(H|_{e^\perp})\ge0.
\tag{7.4}
\]

At a relative face interior, differentiation of the Neumann condition in
every face-tangent direction shows that

\[
H\nu=a\nu,\qquad a=H[\nu,\nu].
\tag{7.5}
\]

Write \(\nu=-ke+su\), where \(s=\sqrt{1-k^2}\) and \(u\perp e\) is a
unit vector. Taking the \(e\) and \(u\) components of (7.5), and
eliminating \(H[e,u]\), gives

\[
(1-k^2)H[u,u]=(1-2k^2)a+k^2q.
\tag{7.6}
\]

The left side is nonpositive, while \(q\ge0\) and \(1-2k^2>0\). Hence
\(a\le0\), and therefore

\[
\partial_\nu h=H[\nu,e]=-ka\ge0
\tag{7.7}
\]

on each relative face interior.

We now pass this face inequality across the polyhedral skeleton. On each
bounded set, the union of strata of codimension at least two has zero
\(H^1\) capacity. Here is the needed explicit approximation. For an
affine subspace of codimension two, take a cutoff equal to zero at
distance at most \(\varepsilon^2\), equal to one at distance at least
\(\varepsilon\), and affine in the logarithm of distance between them.
On a bounded set, its squared-gradient integral is
\(O(1/|\log\varepsilon|)\), while its difference from one tends to zero
in \(L^2\). A finite product treats all the lower-dimensional strata,
each of which lies in such a subspace. Smooth approximations can be
chosen with values between zero and one.

If \(\varphi\) is a nonnegative smooth compactly supported test function
on \(\overline\Gamma\) whose support avoids the skeleton, harmonicity
and (7.7) give

\[
\int_\Gamma\nabla h\cdot\nabla\varphi\ge0.
\tag{7.8}
\]

Multiplying an arbitrary such \(\varphi\) by the preceding cutoffs gives
nonnegative tests converging to it in \(H^1\). Since \(h\) is locally
\(H^1\), (7.8) follows for the original test as well. In particular, no
unaccounted edge flux is discarded.

Use (7.8) with

\[
\varphi_R(x)=e^{-e\cdot x}\chi_0(|x|/R),
\tag{7.9}
\]

where \(\chi_0\) is a smooth nonnegative cutoff equal to one on
\([0,1]\) and zero on \([2,\infty)\). Homogeneity and the local \(H^2\)
estimate give polynomial growth of the relevant annular Sobolev norms.
The first condition in (7.1) gives exponential decay of the weight. The
term containing \(\nabla\chi_0(|x|/R)\) therefore tends to zero by
Cauchy-\/-Schwarz. The remaining term converges, so (7.8) yields

\[
0\le-\int_\Gamma e^{-e\cdot x}\partial_e h
=-\int_\Gamma e^{-e\cdot x}q\le0.
\tag{7.10}
\]

It follows that \(q=0\) in the interior. The negative-semidefinite
transverse Hessian block then has zero trace and must vanish. Writing
\(x=te+z\), \(z\perp e\), connectedness of each transverse slice and
homogeneity give constants \(A\in\mathbb R\) and \(b\in e^\perp\) such
that

\[
\psi(te+z)=At^\beta+t^{\beta-1}b\cdot z.
\tag{7.11}
\]

Harmonicity becomes

\[
0=(\beta-1)t^{\beta-3}
\big(\beta At+(\beta-2)b\cdot z\big).
\tag{7.12}
\]

Since \(1<\beta<2\) and \(\Gamma\) has nonempty interior, \(A=0\) and
\(b=0\). This contradicts the nonzero mode and proves the theorem. \(\square\)

\subsubsection{7.1. Consequence for the small-parameter Robin
question}\label{consequence-for-the-small-parameter-robin-question}

For a consistent-normal cone, let \(\gamma\) satisfy
\(\gamma\cdot\nu_i=-1\) for every face and set \(e=\gamma/|\gamma|\).
The second condition in (7.1) becomes \(|\gamma|>\sqrt2\); the first
condition is a separate strict-dual requirement.

\Needspace{6\baselineskip}\textbf{Corollary 7.2.} Let \(\Omega\subset\mathbb R^d\), \(d\ge3\), be
a bounded convex polyhedron. Suppose the first Neumann variation \(v\)
of its normalized Robin ground state has at a boundary point \(x_0\) a
nonzero leading nonlinear homogeneous Neumann mode \(\psi\) of degree
\(\beta\in(1,2)\). Suppose also that the tangent cone has consistent
normals and satisfies (7.1) with \(e=\gamma/|\gamma|\). Then the Robin
ground state on \(\Omega\) has a nonconvex superlevel set for every
sufficiently small positive Robin parameter.

\emph{Proof.} The expansion and interior derivative estimates in {[}ACH,
Section 9{]} give, on each fixed compact subset of the tangent cone,

\[
\begin{aligned}
\nabla v(x_0+\rho x)&=\gamma+\rho^{\beta-1}\nabla\psi(x)
+o(\rho^{\beta-1}),\\
D^2v(x_0+\rho x)&=\rho^{\beta-2}D^2\psi(x)
+o(\rho^{\beta-2}).
\end{aligned}
\tag{7.13}
\]

The lower-order quadratic particular solution is included in the
remainders. By Theorem 7.1 there are an interior point \(x\) and a
vector \(\xi\perp e\) with \(D^2\psi(x)[\xi,\xi]>0\). Correct \(\xi\) to
\(\xi+a_\rho e\) so that it is orthogonal to \(\nabla v(x_0+\rho x)\).
Since \(\gamma\cdot e=|\gamma|>0\), the first line of (7.13) gives
\(a_\rho=O(\rho^{\beta-1})\). The second line gives strictly positive
second derivative in this corrected direction for small \(\rho\).
Taylor\textquotesingle s theorem at one fixed such interior point
produces two endpoints whose \(v\) values strictly exceed that of their
midpoint.

The Neumann perturbation \(u_\alpha=1+\alpha v+o(\alpha)\) from {[}ACH,
Proposition 3.1{]} holds uniformly on these three fixed points. Thus the
strict midpoint inequality persists for \(u_\alpha\) for all
sufficiently small \(\alpha>0\). This is the transfer mechanism of
{[}ACH, Lemma 9.5{]}. \(\square\)

The leading nonzero mode in this corollary is a hypothesis. No assertion
is made that a chosen vertex necessarily has a nonzero coefficient, or
that every polyhedron has a vertex satisfying (7.1).

\subsubsection{7.2. Nonempty scope and the remaining angle
range}\label{nonempty-scope-and-the-remaining-angle-range}

Theorem 7.1 is not vacuous in its degree interval. Fix \(0<k<1/\sqrt2\)
and put \(s=\sqrt{1-k^2}\). In coordinates \((x,y,t)\), consider the
cone with three outward normals

\[
\begin{aligned}
\nu_1&=(s,0,-k),\\
\nu_2&=(-s\cos\varepsilon,s\sin\varepsilon,-k),\\
\nu_3&=(-s\cos(2\varepsilon),-s\sin(2\varepsilon),-k),
\end{aligned}
\qquad 0<\varepsilon<\pi/6.
\tag{7.14}
\]

The transverse normals positively span \(\mathbb R^2\), so the section
\(t=1\) is a bounded triangle and \(t\) is strictly positive on every
nonzero direction in the closed cone. Thus (7.1) holds with
\(e=(0,0,1)\) and a positive \(\eta\) depending on \(\varepsilon\).

As \(\varepsilon\downarrow0\), the spherical sections converge almost
everywhere to the spherical section of
\(\{|x|<(k/s)t\}\times\mathbb R_y\). On this limiting section, \(y\) has
mean zero and Neumann eigenvalue two. Indeed, latitude integration gives
\(\int y^2=|A_0|/3\), while \(|\nabla_{S^2}y|^2=1-y^2\). The Rayleigh
quotient on the finite sections, tested with \(y\) minus its mean,
therefore tends to two by bounded convergence. On every finite pointed
cone, {[}ACH, Theorem 9.1{]} gives \(\lambda_1>2\), since equality
requires a linear factor. Hence, for sufficiently small positive
\(\varepsilon\),

\[
2<\lambda_1<6,\qquad
1<\beta_1<2,\qquad \beta_1(\beta_1+1)=\lambda_1.
\tag{7.15}
\]

No explicit cutoff in \(\varepsilon\) or eigenfunction parity is needed.

The angle condition cannot simply be deleted from the face-sign step.
For example, with \(e=(1,0,0)\),

\[
\nu=(-\sqrt3/2,1/2,0),\qquad
H=\begin{pmatrix}1&-\sqrt3&0\\-\sqrt3&-1&0\\0&0&0\end{pmatrix},
\tag{7.16}
\]

we have \(\operatorname{tr}H=0\), \(H|_{e^\perp}\preceq0\), and
\(H\nu=2\nu\). The normal second derivative is positive and the desired
face-flux sign reverses. This matrix is an obstruction to that algebraic
proof step, not a counterexample to the cone conjecture. The borderline
angle, the larger-angle range, and cones without the strict-dual
property require further arguments. The general {[}ACH, Section 10,
Conjecture 2{]} remains unresolved here.

\sloppy\section{References}\label{references}

{[}ACH{]} B. Andrews, J. Clutterbuck and D. Hauer, \emph{Non-concavity
of the Robin ground state}, Cambridge Journal of Mathematics \textbf{8}
(2020), no. 2, 243-310.
\href{https://doi.org/10.4310/CJM.2020.v8.n2.a1}{Final journal article}.
\href{https://arxiv.org/abs/1711.02779v2}{arXiv:1711.02779v2}.

{[}CF{]} G. Crasta and I. Fragalà, \emph{Concavity properties of
solutions to Robin problems}, Cambridge Journal of Mathematics
\textbf{9} (2021), no. 1, 177-212.
\href{https://doi.org/10.4310/CJM.2021.v9.n1.a3}{Final journal article}.
\href{https://arxiv.org/abs/2006.07192v1}{arXiv:2006.07192v1}.

{[}D{]} M. Dauge, \emph{Neumann and mixed problems on curvilinear
polyhedra}, Integral Equations and Operator Theory \textbf{15} (1992),
227-261.
\href{https://dauge.pages.math.cnrs.fr/publis/DaugeMixed92.pdf}{Author-posted
text}, Theorem 8.1.

{[}YZ{]} D. Ye and D. Zhang, \emph{Concavity Properties of Robin
Solutions on \(C^{3,1}\) Uniformly Convex Domains},
\href{https://arxiv.org/abs/2609.06223v1}{arXiv:2609.06223v1}, September
5, 2026.

{[}T{]} P. Tolksdorf, \emph{The Stokes resolvent problem: optimal
pressure estimates and remarks on resolvent estimates in convex
domains}, Calculus of Variations and Partial Differential Equations
\textbf{59} (2020), article 154.
\href{https://doi.org/10.1007/s00526-020-01811-8}{Final journal
article}, Section 4.

\end{document}
